
\documentclass[../../../physics-standard_answers_all.tex]{subfiles}

\begin{document}

\setlength{\abovedisplayskip}{2mm}
\setlength{\belowdisplayskip}{1mm}

\begin{enumerate}

    \item Yについての電荷保存則より，
        %
        \begin{align*}
            -x_1 + y_1 = 0 \;.
        \end{align*}
        %
        キルヒホッフ則より，
        %  
        \begin{align*}
            \frac{y_1}{\varepsilon_0 S} \cdot 2d
            + \frac{x_1}{\varepsilon_0 S} \cdot d = 2V_0
            \qquad \text{∴} \; x_1 = y_1
            = \uwave{\frac{2\varepsilon_0 S}{3d} V_0} \;.
        \end{align*}

    \item 電荷保存則より，Xの電荷は変わらない．つまり，
        %
        \begin{align*}
            x_2 = x_1 = \uwave{\frac{2\varepsilon_0 S}{3d} V_0} \;.
        \end{align*}
        %
        キルヒホッフ則より，
        %
        \begin{align*}
            \frac{y_2}{\varepsilon_0 S} \cdot 2d = V_0
            \qquad \text{∴} \; y_2 = \uwave{\frac{\varepsilon_0 S}{2d} V_0} \;.
        \end{align*}

    \item Yについての電荷保存則より，
        %
        \begin{align*}
            -x_3 + y_3 = -x_2 + y_2
            \qquad \text{∴} \; -x_3 + y_3 = -\frac{\varepsilon_0 S}{6d}V_0 \;.
        \end{align*}
        %
        キルヒホッフ則より，
        %
        \begin{align*}
            & \frac{y_3}{\varepsilon_0 S} \cdot 2d
            + \frac{x_3}{\varepsilon_0 S} \cdot d = 2V_0 \\
            & \text{∴} \; x_3 = \uwave{\frac{7\varepsilon_0 S}{9d} V_0} \;,
            \quad y_3 = \uwave{\frac{11\varepsilon_0 S}{18d} V_0} \;.
        \end{align*}

    \item Yについての電荷保存則より，
        %
        \begin{align*}
            -x_4 + y_4 = -x_3 + y_3
            \qquad \text{∴} \; -x_4 + y_4 = -\frac{\varepsilon_0 S}{6d}V_0 \;.
        \end{align*}
        %
        キルヒホッフ則より，
        %
        \begin{align*}
            & \frac{y_4}{\varepsilon_0 S} \cdot 2d
            + \frac{x_4}{\varepsilon_0 \varepsilon_r S} \cdot d = 2V_0 \\
            & \text{∴} \; x_4 = \uwave{\frac{7\varepsilon_r}{3(2\varepsilon_r + 1)}
            \frac{\varepsilon_0 S}{d} V_0} \;,
            \quad y_4 = \uwave{\frac{12\varepsilon_r - 1}{6(2\varepsilon_r + 1)}
            \frac{\varepsilon_0 S}{d} V_0} \;.
        \end{align*}

\end{enumerate}

\end{document}
